\section{Comparación de los tres paradigmas de aprendizaje}

Antes de entrar de lleno al aprendizaje por refuerzo, es necesario compararlo sistemáticamente con los dos paradigmas de aprendizaje vistos previamente, para entender su posición particular.

\begin{figure}[H]
\centering
\includegraphics[width=0.95\textwidth]{slides-images/slide-07.jpg}
\caption{Comparación de los tres tipos de problemas de aprendizaje: Supervised Learning, Unsupervised Learning y Reinforcement Learning\protect\footnotemark}
\end{figure}
\footnotetext{Diapositivas, página 7.}

\subsection{Supervised Learning (aprendizaje supervisado)}

\begin{itemize}
    \item \textbf{Datos}: pares $(x, y)$, donde $x$ es el dato de entrada y $y$ es la etiqueta
    \item \textbf{Objetivo}: aprender una función que mapee $x$ a $y$
    \item \textbf{Analogía}: mostrarle al modelo muchas imágenes de manzanas y decirle «esto es una manzana», de modo que el modelo aprenda a reconocer imágenes nuevas de manzanas
\end{itemize}

\subsection{Unsupervised Learning (aprendizaje no supervisado)}

\begin{itemize}
    \item \textbf{Datos}: solo $x$, sin etiquetas
    \item \textbf{Objetivo}: aprender la estructura y la distribución intrínsecas de los datos
    \item \textbf{Analogía}: mostrarle al modelo muchas imágenes de objetos similares (sin decirle qué son), de modo que el modelo aprenda a descubrir sus características comunes
\end{itemize}

\subsection{Reinforcement Learning (aprendizaje por refuerzo)}

\begin{itemize}
    \item \textbf{Datos}: State-Action pairs (pares de estado y acción)
    \item \textbf{Objetivo}: maximizar la \textbf{recompensa acumulada futura} a lo largo de múltiples pasos temporales
    \item \textbf{Analogía}: no es necesario decirle al modelo «esto es una manzana»; en vez de eso, se le deja aprender mediante la interacción que «comer una manzana proporciona nutrientes y así sobrevive»
\end{itemize}

\begin{importantbox}{Diferencia esencial entre los tres paradigmas}
\begin{itemize}
    \item \textbf{Supervised Learning}: aprende el mapeo de entrada a salida a partir de \textbf{datos etiquetados}
    \item \textbf{Unsupervised Learning}: descubre estructuras ocultas a partir de \textbf{datos sin etiquetar}
    \item \textbf{Reinforcement Learning}: aprende una política de decisión óptima mediante la \textbf{interacción con el entorno}, con el objetivo de maximizar el retorno acumulado a largo plazo
\end{itemize}
El algoritmo de aprendizaje (Algorithm) es independiente (agnostic) de la arquitectura de la red (Architecture). Una misma arquitectura (como CNN) puede usarse en distintos paradigmas de aprendizaje.
\end{importantbox}

\subsection{Resumen de la sección}
Los tres paradigmas de aprendizaje corresponden a distintas formas de datos y objetivos de aprendizaje. Lo particular del aprendizaje por refuerzo radica en que: (1) los datos no se recopilan de antemano, sino que se generan dinámicamente mediante la interacción; (2) el objetivo no es predecir etiquetas ni descubrir estructuras, sino aprender una política de conducta óptima que maximice el retorno a largo plazo; (3) la decisión actual afecta los datos que se observarán después (no son i.i.d.).


